\documentclass{article}
\usepackage[margin=1in]{geometry}
\usepackage{problemset}
\showsolntrue

% Custom commands
\newcommand{\R}{\mathbb{R}}
\newcommand{\X}{\mathcal{X}}
\newcommand{\Y}{\mathcal{Y}}
\newcommand{\E}{\mathbb{E}}
\newcommand{\V}{\text{Var}}
\renewcommand{\P}{\mathbb{P}}
\newcommand{\x}{\mathbf{x}}
\newcommand{\w}{\mathbf{w}}

\begin{document}

\pstitle{%
    DSC 180A%
}{
    Participation 3%
}[due Sunday, October 18 at 11:59PM]

\begin{probset}
% --- Problem 1 ---
\begin{prob}
    Describe one thing you found interesting in the reading. Describe what it is in your own words and why you found it interesting.

    \begin{soln}
        One thing that I found interesting is that the obstacle to low-bit quantization isn't the weights but rather the activations. Each paper handles this by moving the difficulty from activations to weights. I'm not sure I fully understand the exact mechanisms though. As I understood it, all of the papers that we were assigned also attempted to tackle the same problem in different ways, mostly by scaling channels in specific ways.

        This was my first real look into quantization and I had heard about it before, but I plan on continuing to look into this space now.
    \end{soln}
\end{prob}

% --- Problem 2 ---
\begin{prob}
    Describe one thing that you found difficult to understand. Try to be specific about what you don't think you understand.

    \begin{soln}
        The difficulty that I had with these readings was understanding the exact mechanism in which the problem was solved, in this case quantization. I understand the motivation and the problem, but I'm not sure why each specific solution works and wouldn't be able to explain it to someone with my current understanding. 
        
        For example, GPTQ compensates for the error it makes when it rounds a weight by nudging the weights it hasn't quantized yet, and it uses the Hessian to decide how. I follow that the Hessian is measuring sensitivity, but I can't really explain why using its inverse gives the right adjustment, or why it's okay to quantize the columns in a fixed order.
    \end{soln}
\end{prob}

% --- Problem 3 ---
\begin{prob}
    In the reading, identify one choice that the author made in their investigation and your best guess at why they made that choice. For example, this may be a choice in sub-population to study, or a choice in parameter for a model.

    \begin{soln}
        A choice that stood out to me was SpinQuant's decision to make the rotation matrices learnable instead of just fixing them like QuaRot does. QuaRot already showed that a fixed rotation smooths out the outliers and gets you down to 4-bit, so SpinQuant could have stopped there, but they instead trained the rotation on a small calibration set to minimize the quantized model's loss.

        My guess at why they made this choice is that they noticed not every rotation works equally well, so leaving it fixed is leaving accuracy on the table. Since a rotation is reversible no matter what values it ends up with, they can learn a better one without any extra cost at inference time. 
    \end{soln}
\end{prob}
\end{probset}
\end{document}